% ID: SEL-D0012-BA-E02
% ESTADO: validada. S013; tres pasadas del mismo autor, sin segundo revisor.
\textbf{a) Potencias y fórmula general.}\par\nopagebreak[4]
Multiplicando filas por columnas se obtiene
\[
A^2=\begin{pmatrix}2&0&2\\0&1&0\\2&0&2\end{pmatrix},\qquad
A^3=\begin{pmatrix}4&0&4\\0&1&0\\4&0&4\end{pmatrix},\qquad
A^4=\begin{pmatrix}8&0&8\\0&1&0\\8&0&8\end{pmatrix}.
\]
El bloque de las filas y columnas primera y tercera duplica sus entradas en cada multiplicación; la entrada central sigue siendo $1$. Para $n\geq1$,
\[
\boxed{A^n=\begin{pmatrix}2^{n-1}&0&2^{n-1}\\0&1&0\\2^{n-1}&0&2^{n-1}\end{pmatrix}.}
\]
La base $n=1$ reproduce $A$; multiplicar la fórmula por $A$ duplica las cuatro entradas y deja fija la central, lo que prueba la inducción. Si la convención de naturales incluye $0$, se añade por separado $A^0=I_3$; la fórmula anterior no se aplica a $0$.

\textbf{b) Inversa de B.}\par\nopagebreak[4]
Desarrollamos por la primera fila:
\[
\det B=1\begin{vmatrix}1&-1\\1&0\end{vmatrix}
+2\begin{vmatrix}1&1\\2&1\end{vmatrix}=1+2(-1)=-1\ne0.
\]
Por tanto, $B$ es invertible. No es necesario calcular toda su inversa para justificarlo.

\textbf{c) Existencia y resolución.}\par\nopagebreak[4]
$B$ es $3\times3$ y $C$ una columna, así que $X=(u,v,w)^t$ tiene tamaño $3\times1$. Al ser $B$ invertible, existe solución única. Las ecuaciones son
\[
u+2w=1,\qquad u+v-w=-3,\qquad2u+v=1.
\]
De la primera $u=1-2w$ y de la tercera $v=-1+4w$. En la segunda queda $w=-3$; entonces $u=7$ y $v=-13$:
\[
\boxed{X=\begin{pmatrix}7\\-13\\-3\end{pmatrix}.}
\]
Comprobación directa: $BX=(7-6,\ 7-13+3,\ 14-13)^t=(1,-3,1)^t=C$.
